\documentclass{xdyydoc}

\newcommand{\DocDate}{2026-07-30}
\newcommand{\DocVersion}{v0.3.4}
\usepackage{xeCJKfntef, xpinyin}
\usepackage{graphicx}
\usepackage{zhlipsum}
\usepackage{tabularray}
\usepackage{../exam-zh-choices}
\usepackage{../exam-zh-question}
\usepackage{../exam-zh-symbols}
\usepackage{../exam-zh-chinese-english}
\usepackage{../exam-zh-textfigure}
\usepackage{../exam-zh-math}

\ExplSyntaxOn
\NewDocumentCommand \examsetup { m }
  { \keys_set:nn { exam-zh } {#1} }
\ExplSyntaxOff
% \usepackage[
%   backend = biber,
%   style = gb7714-2015
% ]{biblatex}
% \addbibresource{exam-zh.bib}
\graphicspath{{figures}}

\hypersetup{
  pdftitle  = {exam-zh: 中国试卷 LaTeX 模板},
  pdfauthor = {夏康玮}
}
% 全角标点放在引号中，需要改成半角式，否则间距过大，不好看
\def\FSID{“{\xeCJKsetup{PunctStyle=banjiao}。}”} % U+3002
\def\FSFW{“{\xeCJKsetup{PunctStyle=banjiao}．}”} % U+FF0E
\def\COFW{“{\xeCJKsetup{PunctStyle=banjiao}：}”} % U+FF1A
\def\SCFW{“{\xeCJKsetup{PunctStyle=banjiao}；}”} % U+FF1B


\title{\textcolor{MaterialIndigo800}{%
  \textbf{exam-zh: 中国试卷 \LaTeX \xpinyin[font=\sffamily,format=\color{MaterialIndigo800}]{模}{mu2}板}}}
\author{%
  夏康玮\thanks{%
    李泽平构建了 \cls{exam-zh} 的最初的基本框架；张庭瑄开发 \cls{exam-zh-font} 模块；郭李军开发了连线题环境}
}
\date{\DocDate\quad \DocVersion%
  \thanks{%
    GitHub: \url{https://github.com/xkwxdyy/exam-zh} \\
    Gitee（国内镜像）: \url{https://gitee.com/xkwxdyy/exam-zh} \\
    \hspace*{1.5em} QQ 用户交流群：652500180
  }
}

\ExplSyntaxOn
\NewDocumentCommand { \scoringbox } { s }
  {
    \IfBooleanTF {#1}
      { \__examzh_scoringbox_onecolumn: }
      { \__examzh_scoringbox_twocolumn: }
  }
\cs_new_protected:Nn \__examzh_scoringbox_twocolumn:
  {
    \begin{tabular}{|c|c|}
      \hline 
      得分 & \rule{3em}{0pt}\rule[-0.7em]{0pt}{2em} \\\hline
      阅卷人 & \rule{3em}{0pt}\rule[-0.7em]{0pt}{2em} \\\hline
    \end{tabular}
  }
\cs_new_protected:Nn \__examzh_scoringbox_onecolumn:
  {
    \begin{tabular}{|c|}
      \hline 
      得分\rule[-0.7em]{0pt}{2em} \\\hline
      \rule[-0.7em]{0pt}{2em} \\\hline
    \end{tabular}
  }
\ExplSyntaxOff

\AddToHook{env/latexexample/after}
  {%
    \examsetup{
      question/index=1
    }
  }
\usepackage{amssymb}

\begin{document}

% 封面的页边距

\newgeometry{
  left   = 2.2 in,
  right  = 1.25 in,
  top    = 1.25 in,
  bottom = 1.00 in
}

\maketitle

\input{./body/cover.tex}


% 用户手册的页边距

\newgeometry{
  left   = 1.75 in,
  right  = 0.80 in,
  top    = 1.25 in,
  bottom = 1.00 in
}

\tableofcontents


% 介绍
\input{./body/introduction.tex}
% 安装
\input{./body/installation.tex}

% \clearpage
% 使用
\input{./body/usage.tex}


% 宏包依赖情况
\input{./back/package.tex}
% 主要更新
\input{./back/main-changelog.tex}

% 参与开发
\input{./back/development.tex}


% 关于提问
\input{./back/ask-questions.tex}


% 关于作者
\input{./back/about-author.tex}



% \section{\cls{exam-zh}：TODO}

% \begin{itemize}
%   % \item 增加试卷范例（语文，英语）
%   % \item 答案控制
%   %   \begin{itemize}
%   %     \item 选择题
%   %       \begin{itemize}
%   %         \item 题目下方
%   %         \item 括号内
%   %         \item 最后：列表形式、表格形式
%   %       \end{itemize}
%   %     \item 填空题
%   %       \begin{itemize}
%   %         \item 题目下方
%   %         \item 划线内
%   %         \item 最后
%   %       \end{itemize}
%   %     \item 解答题
%   %       \begin{itemize}
%   %         \item 题目下方
%   %         \item 移动到最后
%   %       \end{itemize}
%   %   \end{itemize}
%   \item 选择题答案标记
%   % \item \env{choices} 环境“均分”效果
%   % \item 图文排版（参考 xkwxdyy 的 \pkg{text-figure} 宏包和 qinglee 的 \pkg{wrapstuff} 宏包）
%   \item 测试兼容性
%   \item \env{question} 环境的引用
%   \item 每个子包设置用户接口
%   \item 通过 \tn{CJKunderline} 的 \meta{hidden} 键实现 \tn{fillin} 的无答案的时候，长度是答案的长度
%   \item \env{poem} 环境的引用字体不是斜体
%   \item 用 meta 优化键值，比如统一的 show-answer
%   \item 设置 \tn{vec} 命令的定制开关。
%   \item 学习 randexam 的做法，用 exam-zh 文类来写手册
% \end{itemize}


\end{document}
